% ID: SEL-D0053-OB-E02
% ESTADO: validada. S054.
\textbf{a)} Como $1$ está dentro del tramo cuadrático, el mínimo exige $f'(1)=2-b=0$, luego $b=2$. El coeficiente cuadrático positivo asegura el mínimo en ese tramo. Para la continuidad en $2$, igualamos $5-2b=4+a$, de donde $a=-3$. Con estos valores, $f(1)=0$ y el tramo lineal es creciente a partir del valor $1$ en la unión. Por tanto, el mínimo es también absoluto.

\textbf{b)} En $-2$ se usa la rama $x^2-6x+1$. Su valor es $17$ y su derivada vale $-10$. La ecuación punto-pendiente es $y-17=-10(x+2)$, es decir, $\boxed{y=-10x-3}$.
